% GENERATED FILE. Edit canonical lessons, then run npm run generate:books.
% canonical-source-hash: fnv1a64:e7558e3b4b5d222e
% canonical-lessons: TE-C22-pacha-pasupu, TE-S125-letter-ha

\chapter{Green and Yellow}
\label{ch:te-green-and-yellow}
\hlchaptermodality{22}

\begin{chapteropening}
\textbf{By the end of this chapter:} \emph{I can name green and yellow in Telugu, on top of the four colours I already had.}

Say \te{పచ్చ} and \te{పసుపు} and show they are doublets of one Proto-Dravidian *pac- root — where Kannada's yellow broke away as a Sanskrit loan.
\end{chapteropening}

\section[pacca pasupu]{\te{పచ్చ}, \te{పసుపు} — two colors, one Dravidian root}

\label{lesson:TE-C22-pacha-pasupu}



\begin{quote}
In Kannada, green stayed native but yellow got borrowed from Sanskrit. Telugu tells a tighter story: its green word and its yellow word turn out to be \textbf{doublets of the same root}, wearing two different faces.
\end{quote}

\begin{cousinweb}
\textbf{\te{పచ్చ}} (\textbf{pacca}, \textquotedblleft{}\textbf{green}\textquotedblright{}) comes from \textbf{Proto-Dravidian} \emph{\textbf{*pac-}} — the same root already behind Tamil's \textbf{\ta{பச்சை}} (\emph{paccai}) and Kannada's \textbf{\ka{ಹಸಿರು}} (\emph{hasiru}, via Old Kannada \emph{pasir}). No surprise here; this is the same pan-Dravidian green family from the last lesson.
\end{cousinweb}

\begin{cousinweb}
\textbf{\te{పసుపు}} (\textbf{pasupu}, \textquotedblleft{}\textbf{yellow, turmeric}\textquotedblright{}) is, remarkably, \textbf{also} a direct descendant of that \textbf{same} Proto-Dravidian \emph{\textbf{*pac-}} root — a true \textbf{doublet} of \emph{pacca} itself, not a separate borrowing the way Kannada's \emph{haḷadi} was. Tamil even keeps a matching cognate, \textbf{\ta{பசுப்பு}} (\emph{pacuppu}), sitting alongside its own \emph{paccai}. So where Kannada split green and yellow into two unrelated families (native Dravidian vs. a Sanskrit turmeric-loan), Telugu kept \textbf{both} colors inside the same native word-family — \emph{pacca} and \emph{pasupu} are, at root, doublets of the same word, one settling on \textquotedblleft{}green\textquotedblright{} and the other on \textquotedblleft{}yellow, turmeric-colored.\textquotedblright{}
\end{cousinweb}

\subsection*{Your turn}
Say these aloud:

\begin{itemize}
\raggedright
  \item \textquotedblleft{}\te{పచ్చ}\textquotedblright{} — green, from Proto-Dravidian \emph{pac-}
  \item \textquotedblleft{}\te{పసుపు}\textquotedblright{} — yellow/turmeric, a true doublet of pacca, same root
  \item the contrast with Kannada — Telugu kept both colors in one native family; Kannada's yellow broke away as a Sanskrit loan
\end{itemize}

\begin{tcolorbox}[breakable,colback=teal!4,colframe=teal!35!black,title={Before you move on}]
Are \textbf{\te{పచ్చ}} and \textbf{\te{పసుపు}} two unrelated words, or doublets of the same root? (\textbf{Doublets} — both from Proto-Dravidian \emph{*pac-}.) Does Telugu's yellow word follow the same pattern as Kannada's \emph{haḷadi} (a Sanskrit loan)? (\textbf{No} — Telugu kept its yellow word native, unlike Kannada.) What Tamil word is \emph{pasupu}'s direct cognate? (\textbf{\ta{பசுப்பு}}, \emph{pacuppu}.)
\end{tcolorbox}
\section[ha]{\te{హ} — one character, met inside words you already say}

\label{lesson:TE-S125-letter-ha}



\begin{quote}
Before the new one: \te{ఎ} — what does it do?

One character this time — and you have been saying it for pages without knowing which mark on the page it was.
\end{quote}

\subsection*{Script you'll notice: \te{హ}}
\textbf{\te{హ}} — \emph{ha}.

It is a \textbf{consonant}, and in this script a consonant is never bare: it comes with an \emph{a} already in it. So it is not \emph{h}, it is \textbf{ha}.

You already say these, and both have \te{హ} somewhere inside them:

\begin{itemize}
\raggedright
  \item \textbf{\te{సహాయం}} \emph{sahāyam} — help
  \item \textbf{\te{హైదరాబాద్}} \emph{Hyderābād} — the city from the sentence about where you live
\end{itemize}

Look at where those two came from. \emph{Sahāyam} is Sanskrit; \emph{Hyderābād} reached Telugu through Persian and Urdu. In the words this book has given you so far, \te{హ} sits on borrowings, and the word for noon — the one you say every time you greet somebody in the middle of the day — carries it for the same reason.

\subsection*{Writing: \te{హ} — copy what you see}
Put your pen on \te{హ} and follow its line. Copy the shape you can see — slowly, and larger than it is printed.

\begin{quote}
This book does not yet tell you \textbf{where to start the character or which way to travel}. That is a real thing, taught with real variation from school to school, and it is not written down here until it can be written down with a source. Copying what is in front of you needs no such source, and it is how the shape gets into your hand in the meantime.
\end{quote}

\subsection*{Your turn}
\begin{itemize}
\raggedright
  \item \emph{Look:} at these words, and find \te{హ} in the ones that have it
\end{itemize}

\begin{quote}
\te{సహాయం}  ·  \te{నమస్కారం}
\end{quote}

\begin{itemize}
\raggedright
  \item \emph{Trace it:} \te{హ} three times, saying \emph{ha} as you finish each one
  \item \emph{Look:} back at any page of this chapter and find \te{హ} once more
\end{itemize}

\begin{tcolorbox}[breakable,colback=teal!4,colframe=teal!35!black,title={Before you move on}]
Which character is this — \te{హ}? What sound does it carry? (\emph{\textbf{ha}}.) Name one word you already say that contains it.
\end{tcolorbox}
